%!TEX root = ../main/main.tex

En esta pregunta usted deberá mostrar la correctitud del protocolo RSA entre múltiples partes y luego programar una clase que represente a un participante de dicho protocolo. El objetivo de este protocolo es permitir a un grupo de participantes comunicarse de tal forma que todo participante del grupo puede encriptar mensajes que puede decriptar cualquier participante del grupo.\\

El protocolo RSA entre múltiples partes funciona de la siguiente forma: Supongamos que tenemos $\ell$ participantes. Cada participante $i\in\{1,\ldots,\ell\}$ genera dos primos grandes $P_i$ y $Q_i$ al azar y comparte su \emph{módulo} $N_i=P_i\cdot Q_i$. Luego el grupo verifica que todos los $N_i$ sean primos relativos; de lo contrario el protocolo comienza de nuevo. Finalmente, el grupo se pone de acuerdo en un \emph{exponente público} $e\in\{1,\ldots,\min(N_1,\ldots,N_\ell)\}$ (se puede sacar al azar) y verifican que sea primo relativo con $\phi(N_i)$ para cada $i$. Si no lo es, el protocolo comienza de nuevo. Definimos la \emph{llave pública} para encriptar como $(e,N_1,N_2,\ldots,N_\ell)$.

Si un participante $i\in\{1,\ldots,\ell\}$ quiere encriptar un mensaje $m\in\{0,\ldots,\min(N_1,\ldots,N_\ell)-1\}$, primero debe calcular, para cada $j\in\{1,\ldots,\ell\}$, el valor $N_{-j}$ que es la multiplicación de todos los valores $N_k$ con $k\in\{1,\ldots,\ell\}\setminus\{j\}$. Luego, computa $M_j=N_{-j}^{-1}\!\!\mod N_j$. Finalmente, envía el mensaje 
$$c=m^e\cdot M_1\cdot N_{-1} + m^e\cdot M_2\cdot N_{-2} + \cdots + m^e\cdot M_\ell\cdot N_{-\ell}$$

\begin{enumerate}
	\item \text{[2 puntos]} Explique cómo puede un participante decriptar el mensaje $c$. Explique cómo se vería su llave secreta y demuestre que el protocolo es correcto (es decir, al decriptar se obtiene el mensaje original). Es importante que su protocolo pueda ser implementado con la interfaz que se describe en la pregunta (b).
        %Aquí sí puede utilizar el teorema chino de los restos.
	\item \text{[4 puntos]} Escriba una clase en Python que represente a un participante de este protocolo. En particular, su clase deberá respetar al menos la siguiente interfaz.
	\begin{python}
	class MultiPartyRSAParticipant:
		def __init__(self, bit_length: int) -> None:
			# Initializes P, Q and N such that N has at least bit_length and at most bit_length + 1 bits

		def get_N(self):
			# Returns N

		def join_group(self, public_exponent: int, group_moduli: List[int]) -> boolean:
      # Raises an AssertionError if a group has already been joined.
			# Verifies that N occurs exactly once in group_moduli and that all other moduli are coprime with N. Otherwise, returns False.
			# Verifies that e is coprime with (P - 1)(Q - 1). Otherwise, returns False.
			# Precomputes and stores all relevant information for encryption and decryption, then returns True.

		def enc(self, m: int) -> int:
			# Encrypt m according to the multiparty RSA protocol. If no group has been joined raise an AssertionError.
		
		def dec(self, c: int) -> int:
			# Decrypt c according to the multiparty RSA protocol. If no group has been joined raise an AssertionError.
	\end{python}


	A continuación se muestra un bloque de código que utiliza la clase anterior, ejemplificando el funcionamiento del protocolo:

	\begin{python}
		import random

		participant_number = 5
		bit_length = 2048

		participants = []
		for i in range(participant_number):
			participants.append(MultiPartyRSAParticipant(bit_length))

		moduli = [p.get_N() for p in participants]
		public_exponent = random.randint(1, min(moduli))

		for p in participants:
			p.join_group(public_exponent, moduli)

		encryptor = random.choice(participants)

		plaintext = random.randint(0, min(moduli) - 1)
		cyphertext = encryptor.enc(plaintext)

		for p in participants:
			if p.dec(cyphertext) != message:
				raise AssertionError("Incorrect decryption")

		print("All good")
		
	\end{python}
\end{enumerate}
